\begin{tcolorbox}[breakable, colback=white, colframe=black!80,
  boxrule=0.6pt, arc=0mm, left=3mm, right=3mm, top=2mm, bottom=2mm]
\ttfamily\small\raggedright
After the mid-{}19th century, large-{}scale Zenpokoenfun began to appear in Western Japan. Examples of Kansai region include Kofun of Tokuchi in Sakurai, Nara Prefecture, the Kofun of Tsui Otsuzuyama in Kizugawa, Kyoto Prefecture, the Kofun of Ujima Tezuyama in Okayama, Okayama Prefecture, and the Kofun of Ishibutai in Kanda, Fukuoka Prefecture. At the same time, there were also Zenpokoenfun in Chūgoku region and Kyushu. For example, the ancient tombs of Motoshi Inawa in the city of Asahi, Kyoto Prefecture, and the ancient tombs of Tomikaga in the city of Okayama, among others, the largest one is the Tokutoma ancient tomb. The total length of the mound is about 280 meters, which is more than three times larger than the largest mound tomb of about 80 meters in the later period of the Yayoi period, and its area and capacity are far greater than the latter. Although the mound tombs with a front and back circular shape were already scattered throughout Japan during the Yayoi period, the early forms of the Zenpokoenfun were mostly vertical cave-{}style stone chambers with bamboo-{}shaped wooden coffins, which were different from the mound tombs during the Yayoi period. Moreover, early Zenpokoenfun had common features, such as a lower and wider front part compared to the back circular part, and the combination and position of grave goods were the same. Inside the coffin were jade and mirrors, while outside the coffin were a large number of triangular edge divine beast mirrors, iron weapons, agricultural and fishing tools, and so on. Regarding this, Taiichiro Shiraishi believed that at that time, various forces in Kino, led by the Yamatai Kingdom, and the Seto Inland Sea formed an alliance, and obtained iron and various cultural relics from Korea by defeating forces such as Nukoku and the Ito Kingdoms, which controlled the Genkai Sea. On the other hand, Yoshiro Kondo also mentioned that most of the Zenpokoenfun built on flat land were well organized, and the mounds were also built layer by layer after consolidation. Considering the slope and arrangement of the mounds, circular mounds were constructed. In addition, bronze mirrors from mainland China and the \textquotedbl{}Wajinden\textquotedbl{} from the same period indicate that there was an exchange between the two countries at that time. Therefore, it is speculated that the early appearance of the Zenpokoenfun may have been influenced by civil engineering technology from mainland China, and the political and ideological influence is reflected in the joint construction of the Zenpokoenfun by various alliance members as proof of the alliance.\par
\medskip
At the same time, there are also tombs 5, 4, and 3 of the Kamen Mun in the front of the Houyuan Tomb in Eastern Japan during the same period as Ishibutai Kofun. Among them, tomb 5 is the earliest Kofun in Eastern Japan\textquotesingle{}s history, indicating a possible connection between the Kazusa Province and the Yamatai Kingdom at that time. In the 4th century, Hokuriku region, Kantō region, Tōkai region and Tōhoku region were mainly composed of front and rear tombs, which gradually became Zenpokoenfun two to three generations later. Regarding this, Yamato Takeru pointed out that although the \textquotedbl{}Kojiki\textquotedbl{} and \textquotedbl{}Nihon Shoki\textquotedbl{} mentioned that the Four Generals and Japanese Takezun and others had merged East Japan into the territory of the Japanese monarchy through multiple expeditions, he believed that it was actually Kununokuni located in the Nōbi Plain that formed an alliance in East Japan mainly consisting of rear graves in the past and that after the death of Himiko, they lost the battle with the Yamatai Kingdom or negotiated peace under the leadership of the Yamato Kingship. In the end, the alliance of the Dog Slave Kingdom and the alliance of the Evil Horse Kingdom merged into one, forming the Japanese monarchy. Then, the leaders of East and West Japan were distinguished by the Zenpokoenfun and the front and rear tombs. From the end of the 4th century to the 5th century, large Zenpokoenfun resembling Kofun in Mount Taida began to appear in present-{}day Gunma Prefecture. Large Zenpokoenfun have also been built in places such as Yamanashi Prefecture, Ibaraki Prefecture, and Chiba Prefecture.\par
\medskip
After the mid-{}4th century, the Dawa Tombs, originally centered around the Yamato Kofun in the southeast of the Nara Basin and the Ryuben Kofun, gradually shifted to the Daisen Kofun in the north. The large Zenpokoenfun of Dawa was the Furuichi Kofun Cluster and Mozu Tombs. The center of gravity of the Great King\textquotesingle{}s Tomb shifted again in the 5th century to the ancient city tomb group and the hundred-{}tongued bird tomb group in the Osaka Plain. Regarding the reasons for the multiple transfers of these large Zenpokoenfun, based on research on the \textquotedbl{}Kushiji\textquotedbl{} and \textquotedbl{}Nihonshu Ji\textquotedbl{}, the Japanese academic community has proposed views such as the theory of dynasty alternation and the theory of equestrian tribes conquering dynasties to explain. On the other hand, Hirose Kazuo proposed that the structures of the Sasaki, Koji, Shizuki, and Makino ancient tomb groups were similar, and from the late 4th century to the second half of the 5th century, the heads of the four ancient tomb groups jointly played a role in the Yamato Kingship. After death, they deified as the guardian deity of the Yamato royal power and also proposed that the construction of the large Zenpokoenfun had a deterrent effect on various regions. At the same time, in the second half of the 5th century, no large-{}scale Zenpokoenfun were constructed in areas such as Kibi Province, Upper Maoye, and Hyūga Province, and the scale of the Zenpokoenfun in Gokishichidō gradually decreased. According to the inscriptions on Ji Ji, the iron sword unearthed from the Inawa Mountain Kofun, and the large sword unearthed from the Eta Funayama Kofun, Emperor Yūryaku at that time regarded himself as the king of the world, and his authority was reflected in the fact that only he continued to build large Zenpokoenfun in the Japanese archipelago.\par
\medskip
In the 6th century, the horizontal cave-{}style stone chamber of the Kinai system began to be popularized in Zenpokoenfun. At the same time, they began to appear in areas such as Settsu Province, Owari Province, Uji Province, and even in Kantō region, which were considered as bases for the establishment of the Emperor Keitai. Among them, Ueno had a particularly large number of Zenpokoenfun, compared to only 39 in Kōzuke Province during the same period. Ueno alone had 97, while in Kanto there were 216. As both Kino and Oizhang decreased, Only the Kanto region continued to construct large Zenpokoenfun throughout the 6th century, indicating that the Yamato monarchy at that time relied heavily on the Eastern Kingdom in both economic and military aspects. From the end of the 6th century to the beginning of the 7th century, large square and circular tombs replaced the front and became the mainstream. In the middle of the 7th century, they evolved into octagonal tombs. Regarding the reason for the cessation of construction of Zenpokoenfun, Taiichiro Shiraishi believes that it is related to the reforms of Crown Prince Shōtoku and Soga no Umako, as well as the establishment of the national manufacturing system, while the final Zenpokoenfun are believed to be Asama Mountain Kofun, and Rongmachi claims that tKofun were built in the first half of the 7th century.\par
\bigskip
Sumber:\par
\hspace*{1.5em}- \url{https://en.wikipedia.org/wiki/Zenp\%C5\%8Dk\%C5\%8Denfun}
\end{tcolorbox}
